\documentclass{report}
\usepackage{amsmath,amssymb}
\setlength{\parindent}{0mm}
\setlength{\parskip}{1em}
\begin{document}
\begin{center}
\rule{6in}{1pt} \
{\large Mark Berrill \\
{\bf Adaptive Magnetohydrodynamics Simulations with SAMRAI}}

Oak Ridge National Laboratory \\ PO BOX 2008 MS6164 \\ Oak Ridge \\ TN 37831-6164
\\
{\tt berrillma@ornl.gob}\\
Luis Chacon\\
Bobby Philip\end{center}

A wide variety of plasma problems require both high resolution and large
spatial domains that present a problem for existing models. Adaptive Mesh
Refinement (AMR) provides an efficient method of obtaining both high
resolution and a large spatial domain. AMR allows for the possibility of
simultaneously achieving high resolution over part of a spatial domain,
while maintaining a large computational domain. For structure AMR, this
is achieved by providing different levels of refinement over different
parts of the computational domain by dividing the domain into different
patches and then providing refinement over part of those patches. To
achieve these goals we are combining a MHD plasma model with structured
adaptive mesh refinement through the SAMRAI framework [1]. This allows
for the continued use of a relatively simple discretization of the
spatial domain with the flexibility of a spatially and time-varying
resolution. It also lends itself to large scale parallelism through
splitting of the domain and spatial scales. The coupled model uses SAMRAI
to provide a structured AMR hierarchy that maintains the global problem
solution. For each patch, we apply the physics operators by using an
existing magneto-hydrodynamic (MHD) model PIXIE3D [2]. PIXIE3D is a fully
implicit, fully nonlinear finite-volume code that solves the extended MHD
set of equations in generalized curvilinear geometry [2,3]. Its
non-linearly stable collocated discretization, combined with its fully
nonlinear character, lends it exceptional numerical robustness and
conservation properties [2]. Algorithmically, PIXIE3D has demonstrated
excellent scalability serially and in parallel with up to 4092 processors
[4]. Global time integrators and implicit solvers over the entire AMR
hierarchy are then applied to solve the resulting problem over the
complete domain.

We test the algorithm using the island coalescence MHD problem [4]. The
island coalescence is a self-contained, multiscale MHD driven
reconnection problem in which two co-flowing current channels attract and
reconnect. In the reconnection process, extremely thin current sheets of
thickness sqrt(eta) (with eta the normalized resistivity) develop
non-linearly. Capturing such nonlinear current sheet thinning behavior is
essential for the accurate modeling of the macroscopic dynamics, and thus
mesh adaptivity is of the essence. Furthermore, because the MHD system is
a stiff hyperbolic system, finer meshes are accompanied by more stringent
CFL time step stability constraints, and hence an implicit temporal
scheme is desirable.

1. http://computation.llnl.gov/casc/SAMRAI/index.html

2. L. Chac�n. Comput. Phys. Comm., 163, p. 143 (2004).

3. L. Chac�n. Phys. Plasmas, 15, 056103 (2008).

4. J. Finn and P. Kaw, Phys. Fluids, Vol. 20, p. 72 (1977).


\end{document}
